%% kapa: exam=1 round=32 year=2021 period=1 form=A subject=real_estate subject_ko=부동산학원론 seq=2 number=82
%% source: Q-Net 감정평가사 기출문제 (cst00309) artlSeq=5208397 file=extracted/1008/5208397/A형/2021년도 제32회 감정평가사 1차 1교시 A형.pdf
%% conversion: kordoc 4.24.51 -> tools/gen1.py (auto); review=proofread-sonnet; flags=none
\begin{question}{82}
토지의 분류 및 용어에 관한 설명으로 옳은 것은?
\begin{choices}
  \item 필지는 법률적 개념으로 다른 토지와 구별되는 가격수준이 비슷한 일단의 토지이다.
  \item 후보지는 부동산의 용도적 지역인 택지지역, 농지지역, 임지지역 상호간에 전환되고 있는 지역의 토지이다.
  \item 나지는 「건축법」에 의한 건폐율·용적률 등의 제한으로 인해 한 필지 내에서 건축하지 않고 비워둔 토지이다.
  \item 표본지는 지가의 공시를 위해 가치형성요인이 같거나 유사하다고 인정되는 일단의 토지 중에서 선정한 토지이다.
  \item 공한지는 특정의 지점을 기준으로 한 택지이용의 최원방권의 토지이다.
\end{choices}
\end{question}
\answer{2}
